\begin{statementbox}
\textbf{\textcolor{CStatement}{条件与符号}}\par\smallskip
\begin{itemize}[leftmargin=*,itemsep=.35em]
\item 抛物线：$y^2=2px$（向右）或 $y^2=-2px$（向左），其中 $p>0$。这里的\textbf{$2p$ 就是方程中 $x$ 项系数的绝对值}。
\item $A,B$ 是同一条\textbf{过焦点 $F$ 的直线}与抛物线的两个不同交点；$L=|AB|$；直线倾斜角满足 $0<\theta<\pi$（不取轴线方向）。
\item 焦点分别为 $F(\frac p2,0)$、$F(-\frac p2,0)$；相应准线分别为 $x=-\frac p2$、$x=\frac p2$。
\end{itemize}
\par\medskip
\textbf{\textcolor{CStatement}{核心结论：焦点弦长公式}}\par
\[
\boxed{\displaystyle L=\frac{2p}{\sin^2\theta}\qquad(0<\theta<\pi).}
\]
\textbf{\textcolor{CStatement}{极值与通径}}\par
\[
0<\sin^2\theta\leq1\quad\Longrightarrow\quad
\boxed{L\geq2p},\qquad L=2p\iff\theta=\frac\pi2.
\]
等号对应\textbf{过焦点且垂直于抛物线对称轴的弦}，称为\textbf{通径}。
\par\medskip
\textbf{\textcolor{CStatement}{常用等价形式及适用条件}}\par
\begin{itemize}[leftmargin=*,itemsep=.2em]
\item 若直线斜率 $k=\tan\theta$ 存在且 $k\ne0$：
$\displaystyle L=\frac{2p(1+k^2)}{k^2}$。\textbf{竖直直线 $x=\pm p/2$ 不用 $k$，直接取 $L=2p$}。
\item 若写作 $x^2=\pm2py$（上下开口），而 $\theta$ \textbf{仍从 $x$ 轴量起}：
$\displaystyle L=\frac{2p}{\cos^2\theta}$，适用 $\theta\ne\pi/2$。
\end{itemize}
\textbf{易记一句：}左右开口看 $\sin^2\theta$；上下开口看 $\cos^2\theta$；\textbf{永远是与对称轴垂直的通径最短}。
\end{statementbox}
